\chapter[METHODOLOGY]{Methodology}
\label{ch:methodology}

\section{Research Design}
\label{sec:research-design}

This research is quantitative and experimental. Its central claim --- that
fusing RGB, thermal and LiDAR measurements improves the detection of corroded
regions and the estimation of their depth --- is a claim about measurable
performance, so it is settled by a controlled comparison rather than by
argument. The design therefore has one spine: a set of models that differ only
in which sensing modalities they receive, trained and evaluated under identical
conditions, with the metrics and the decision rule fixed before any model is
trained.

The three objectives map onto three dependent workstreams. Objective~1 produces
the registered trimodal data without which no comparison is possible.
Objective~2 produces the model that consumes it. Objective~3 is the comparison
itself, and it is the only one of the three whose outcome is genuinely open: a
finding that the third modality does not repay its acquisition and registration
cost is a valid answer to Research Question~4, and the protocol is written so
that such a finding would be reported rather than absorbed.

Two design decisions follow from Chapter~\ref{ch:literature-review} and are
stated here because they constrain everything downstream. First, every modality
must be physically sensed. A channel generated from another channel is not a
second sensor, and the reviewed literature contains enough instances of derived
channels presented as modalities that the distinction has to be made explicit
rather than assumed. Second, fusion is performed at the input where the
inspection-domain evidence supports it: \citet{Lim2021} stack chroma and thermal
planes before any feature extraction, \citet{Yang2023} find early fusion
outperforms later stages on building-envelope defects, and \citet{Maru2025}
project image colour onto LiDAR returns before segmentation. Three independent
inspection-domain studies point the same way, which is a stronger basis for the
choice than the autonomous-driving literature's preference for intermediate
fusion.

\section{Project Activities}
\label{sec:project-activities}

\subsection{Data Acquisition}

No public dataset pairs registered RGB, thermal and LiDAR measurements of
corroded metal. That absence is part of the gap this research addresses, and it
makes acquisition an activity of the research rather than an input to it. It is
also the largest single risk to the schedule, and it is treated as such below.

The primary route is a fixed multi-sensor rig imaging corroded and sound steel
specimens at a controlled standoff. The sensor set is an RGB camera, a long-wave
thermal camera and a LiDAR unit, mounted rigidly so that the extrinsic
relationships are solved once rather than per capture --- the arrangement
\citet{Nooralishahi2022} use for drone-based industrial inspection. Specimens
span artificially corroded coupons, whose reference depths can be measured
directly, and, where access permits, in-service structures, which supply the
surface complexity that laboratory coupons lack. \citet{Lim2021} obtained a
publishable result on two bridges from thirty-six acquired image pairs, which
sets a realistic floor on how much data an initial study requires.

Two acquisition questions are answered here rather than deferred, because both
determine whether the method transfers to field inspection at all.

\textbf{Excitation.} Every reviewed thermographic result on steel corrosion was
obtained under deliberate thermal excitation: halogen lamps at one metre
\citep{Lim2021}, quartz-tube heaters \citep{Elizalde2023}, and aerosol cryogen
\citep{Doshvarpassand2022}. No reviewed study recovers subsurface corrosion in
steel passively. Whether a passive regime is viable at inspection standoff is
therefore an open methodological question this research must answer, not a
capability it may assume. The acquisition protocol runs both regimes on the same
specimens so that the answer is measured.

\textbf{Timing.} Under active excitation the thermal measurement is a sequence
rather than a frame: it runs from the stimulation event until the surface
returns to ambient \citep{Doshvarpassand2022}. The thermal stream is then keyed
to an event that the RGB and LiDAR streams are not. Bringing the three into
correspondence is consequently a question of time as well as of geometry, and
the acquisition protocol must define a common time base or state that the regime
is passive and why.

Should field access prove unobtainable within the schedule, evaluation falls
back to laboratory specimens alone, and the generalisation limitation that
\citet{Nash2022} document for single-site training data is reported as a
constraint on the result rather than discovered after it.

\subsection{Calibration and Registration}

The three sensors see different things, at different resolutions, through
different fields of view, and the standard solution does not transfer. A printed
chessboard has no temperature contrast, so it is invisible to a thermal camera;
the target-based procedure that calibrates an RGB--LiDAR pair therefore cannot
be extended to the third sensor by adding a sensor. This is the substance of
Research Question~1 and the reason Objective~1 exists as an objective rather than
as a preprocessing footnote.

Two routes are evaluated against each other. The first is targetless extrinsic
calibration of exactly this sensor triple, for which \citet{Yang2024} report a
trimodal error budget of 2.77\,cm and 0.21\textdegree between LiDAR
and RGB against 3.86\,cm and 0.46\textdegree between LiDAR and
thermal. The thermal channel is the worse-registered one in that budget, by
roughly forty per cent in translation and a factor of two in rotation, which is
consistent with it being the channel without a conventional target. The second
route is a target that the thermal camera can resolve: the self-heating
reference board of \citet{Nooralishahi2022}, whose pattern is sharp in both
bands. Per-capture residual misalignment is then corrected by the coarse-to-fine
multispectral registration of \citet{Li2024c}, which is one of the few methods
in the reviewed literature that treats mismatched resolution and field of view
as the problem to solve rather than as a given.

Registration accuracy is reported as a result of Objective~1 in its own right,
in pixels for the image pair and in millimetres for the projection of LiDAR
returns onto the image plane. Reporting it matters beyond bookkeeping: the
reviewed literature contains no published relationship between residual
registration error and detection accuracy for this sensor set, so the pairing of
a measured alignment error with a measured detection score is itself a
contribution.

\subsection{Preprocessing and Annotation}

Two annotation products are required, not one. Bounding boxes around corroded
regions supply the detection target. A reference maximum pit depth within each
box supplies the regression target, and it is the harder of the two to obtain,
because it cannot be read off an image.

The reference instrument for depth is named here so that the error figures
reported under Objective~2 are interpretable: an estimate quoted to a tenth of a
millimetre is meaningless against ground truth of unstated uncertainty. The
datum is likewise defined explicitly --- depth is measured from the reconstructed
original surface plane of the specimen, not from the local surface, since the
local surface is what corrosion has removed. Both choices are stated before
measurement, and the instrument's own uncertainty is reported alongside every
depth result.

Preprocessing brings the three streams into the common frame established above
and normalises each channel independently, because their dynamic ranges have
nothing to do with one another. Augmentation that assumes natural-image colour
statistics is disabled across all configurations: it is not a meaningful
transformation on a thermal plane or a LiDAR depth plane, and applying it to
some arms and not others would introduce exactly the confound the ablation
exists to avoid.

\subsection{Model Development}

The detection task is object detection --- a bounding box per corroded region,
single class --- and not semantic segmentation. Box annotation costs a fraction
of pixel masks, which matters when the dataset does not yet exist, and mAP is
the metric the fusion literature reports, so the ablation results can be set
against published work.

Development proceeds from a single-modality baseline to the fused model. The
baseline is a conventional RGB detector, which fixes the level that fusion must
exceed to be worth its cost; \citet{Alvani2026} and \citet{Imran2025} establish
what appearance-only methods achieve on corrosion, and \citet{Imran2025} states
from inside the domain that single-modality RGB is the ceiling his framework
runs into.

The fused model extends that baseline with the thermal and geometric channels at
the input, and carries a second head that regresses one scalar per detection:
the maximum pit depth inside the box. The depth value is regressed from the
fused representation rather than read off the LiDAR returns directly. This is
deliberate. A minimum over the LiDAR points inside a box would be arithmetic,
not a learned quantity; it could not be ablated, and removing the geometric
channel would leave the operator with no operand rather than with a degraded
estimate, which would make the depth arm of Objective~3 vacuous. Because the
scalar is regressed, every configuration in the ablation carries a depth head,
and the RGB-only arm is a real comparator rather than a straw one ---
\citet{Egodawela2025} regress metal-loss depth from appearance alone at
0.2829\,mm, which is the existence proof that appearance-only
depth regression is a comparator worth beating.

\subsection{Ablation Protocol}
\label{sec:ablation-protocol}

This is the methodological contribution of the research, and it is specified in
full before any result exists. Its purpose is to remove the freedom to choose,
after seeing the numbers, which comparison to report.

\textbf{Configurations.} Six arms are trained: RGB only; thermal only; LiDAR
only; RGB with thermal; RGB with LiDAR; and the full trimodal set. The three
single-modality arms are controls rather than candidates. They answer a question
the dual and trimodal arms cannot: whether the first modality already suffices,
which is the condition under which fusion has been observed to add nothing.

\textbf{The blank-channel control.} One further arm is trained in which the
additional channel is present but carries no information. Its purpose is to
separate the value of the information from the value of the machinery: a fusion
architecture has more capacity and more parameters than a single-stream one, and
without this control a gain attributable to capacity would be credited to the
sensor. This control is absent from almost all of the reviewed fusion
literature.

\textbf{Held constant.} Architecture, optimiser, learning-rate schedule, epoch
budget, batch size, image size, augmentation, early-stopping criterion and the
data split. The split is drawn once with a fixed seed and reused unchanged by
every arm. Weights are initialised identically across arms; where pretrained
weights are used they are reported separately, because weights fitted to natural
RGB statistics hand the RGB arm a transfer advantage unrelated to the
information in its channels.

\textbf{Metrics, fixed now.} For detection: mAP at 0.5 and mAP averaged
over 0.5 to 0.95, with precision and recall reported alongside. Both
thresholds are reported rather than the looser one alone, because thermal and
LiDAR are expected to sharpen box placement more than they change whether a
region is found at all, and mAP@0.5 would hide that. For depth: mean
absolute error in millimetres against the named reference instrument, reported
with the instrument's own uncertainty, and reported separately for shallow
defects --- \citet{Maru2025} report 7.8\% mean depth error rising to
17.2\% at 0.5\,mm, so an aggregate figure would
conceal failure in exactly the early-stage case worth catching.

\textbf{Repeats.} Every arm is trained with three seeds. Results are reported as
the mean over seeds with the observed range, never as a single best run.

\textbf{The decision rule.} A modality earns its place if the arm containing it
exceeds the best arm without it, on the primary metric, by more than the largest
between-seed spread observed in any arm. Stating the rule in advance and
comparing the gain against the run-to-run variation the experiment itself
exhibits is what stops a difference smaller than noise from being reported as an
effect. Applied to each added channel in turn, the rule also prevents a channel
being credited with the residual of the one before it.

This is not a formality. \citet{Rashidi2026} find a single modality insufficient
on unseen structures under leave-one-bridge-out validation, while recent
radiometric RGB--thermal work on building facades reports a gain confined to one
defect category with no improvement in aggregate accuracy over RGB alone. The
literature supports both outcomes, which is precisely why the question is worth
asking empirically.

\subsection{Validation and Benchmarking}

Validation is held out at the level of the structure, not the image. Random
image-level splits leak: adjacent frames of the same corroded region appear on
both sides of the split, and the reported score then measures memorisation.
\citet{Rashidi2026} hold out whole bridges for this reason, and their conclusion
--- that single-modality performance does not survive the move to an unseen
structure --- only follows because of how they split.

Benchmarking is against three references. The RGB-only arm of this research's
own ablation is the internal baseline. Published appearance-only corrosion
detection results supply the external one. The fused arm is set against
\citet{Lim2021}, the closest published RGB-plus-thermal corrosion result, with
the caveats of that comparison stated: a different structure type, a different
detector generation, and an active excitation regime.

Robustness is assessed by degrading each channel rather than removing it, which
yields a gradient instead of a binary and locates the point at which a modality
stops contributing. For the thermal channel the degradations are reduced
excitation and increased ambient variation; for the geometric channel, added
range noise and reduced point density; for all channels, controlled
misregistration, which produces the error-versus-accuracy relationship the
reviewed literature does not contain.
\section{Objectives and Key Milestones}
\label{sec:milestones}
% -----------------------------------------------------------------
% The Word template carries a SmartArt diagram mapping three
% objectives onto five milestones. The TikZ figure below reproduces
% that mapping -- edit the labels once the objectives are settled.
% -----------------------------------------------------------------

\begin{figure}[htbp]
\centering
\begin{tikzpicture}[
  obj/.style={draw,thick,rounded corners,align=center,
              minimum width=3.4cm,minimum height=1.15cm,
              fill=blue!8,inner sep=3pt},
  ms/.style={draw,thick,rounded corners,align=center,
             minimum width=5.0cm,minimum height=0.9cm,
             fill=orange!12,inner sep=3pt},
  arr/.style={-{Latex[length=2mm]},thick,gray!65}
]
% --- milestones: an evenly spaced ladder on the right ---
\node[ms] (m1) at (0,0) {\footnotesize\textbf{M1} Dataset assembled \& annotated};
\node[ms,below=0.32cm of m1] (m2) {\footnotesize\textbf{M2} Registration pipeline validated};
\node[ms,below=0.32cm of m2] (m3) {\footnotesize\textbf{M3} Baseline \& fused models trained};
\node[ms,below=0.32cm of m3] (m4) {\footnotesize\textbf{M4} Ablation study completed};
\node[ms,below=0.32cm of m4] (m5) {\footnotesize\textbf{M5} Results validated \& benchmarked};

% --- objectives: centered against the milestones each one owns ---
\node[obj] (o1) at ($(m1.west)!0.5!(m2.west) + (-2.5cm,0)$)
  {\textbf{Objective 1}\\[-2pt]\footnotesize Registration \&\\[-3pt]\footnotesize preprocessing};
\node[obj] (o2) at ($(m3.west) + (-2.5cm,0)$)
  {\textbf{Objective 2}\\[-2pt]\footnotesize Detection and\\[-3pt]\footnotesize depth model};
\node[obj] (o3) at ($(m4.west)!0.5!(m5.west) + (-2.5cm,0)$)
  {\textbf{Objective 3}\\[-2pt]\footnotesize Modality\\[-3pt]\footnotesize ablation};

\draw[arr] (o1.east) -- (m1.west);
\draw[arr] (o1.east) -- (m2.west);
\draw[arr] (o2.east) -- (m3.west);
\draw[arr] (o3.east) -- (m4.west);
\draw[arr] (o3.east) -- (m5.west);
\end{tikzpicture}
\caption{Mapping of research objectives to key milestones.}
\label{fig:objectives-milestones}
\end{figure}

\begin{table}[htbp]
\centering
\caption{Key milestones and expected completion.}
\label{tab:milestones}
\begin{tabularx}{\textwidth}{@{}l X l@{}}
\toprule
\textbf{Milestone} & \textbf{Description} & \textbf{Target} \\
\midrule
M1 & Registered trimodal dataset assembled and annotated & Jul 2027 \\
M2 & Registration and preprocessing pipeline validated & Aug 2027 \\
M3 & Baseline (RGB-only) and fused detection--depth models trained & Nov 2027 \\
M4 & Modality ablation over detection and depth completed & Feb 2028 \\
M5 & Results validated and benchmarked & Mar 2028 \\
\bottomrule
\end{tabularx}
\end{table}

\section{Study Plan}
\label{sec:study-plan}

The research runs from May 2026 to April 2028. Figure~\ref{fig:gantt} sets out
the schedule, and Table~\ref{tab:milestones} lists the milestones that mark its
completion points.

% --- Gantt palette and legend keys -------------------------------
\definecolor{gLit}{HTML}{4E79A7}
\definecolor{gProp}{HTML}{3D8C86}
\definecolor{gProg}{HTML}{B07AA1}
\definecolor{gData}{HTML}{59A14F}
\definecolor{gModel}{HTML}{F28E2B}
\definecolor{gThesis}{HTML}{9C755F}
\definecolor{gPhase}{HTML}{4A4A4A}
\definecolor{gMile}{HTML}{C0392B}
\newcommand{\gkey}[1]{\tikz[baseline=-0.55ex]{\fill[#1,draw=black!55,line width=0.3pt]
  (0,-0.09) rectangle (0.40,0.13);}}
\newcommand{\gdia}{\tikz[baseline=-0.55ex]{\fill[gMile,draw=black!55,line width=0.3pt]
  (0.00,0.02) -- (0.11,0.15) -- (0.22,0.02) -- (0.11,-0.11) -- cycle;}}

\begin{landscape}
\begin{figure}[p]
\centering
\begin{ganttchart}[
  hgrid, vgrid,
  x unit=0.70cm,
  y unit title=0.50cm,
  y unit chart=0.40cm,
  title height=1,
  bar height=0.62,
  group height=0.5,
  milestone height=0.5,
  bar/.append style={draw=black!55,thin},
  group/.append style={fill=gPhase,draw=gPhase},
  milestone/.append style={fill=gMile,draw=black!60},
  title label font=\footnotesize,
  group label font=\bfseries\scriptsize,
  bar label font=\scriptsize,
  milestone label font=\scriptsize\itshape,
  title/.append style={fill=black!8},
  bar label node/.append style={align=left,text width=5.2cm},
  group label node/.append style={align=left,text width=5.2cm},
  milestone label node/.append style={align=left,text width=5.2cm}
]{1}{24}
\gantttitle{2026}{8}
\gantttitle{2027}{12}
\gantttitle{2028}{4} \\
\gantttitle{M}{1} \gantttitle{J}{1} \gantttitle{J}{1} \gantttitle{A}{1} \gantttitle{S}{1} \gantttitle{O}{1} \gantttitle{N}{1} \gantttitle{D}{1} \gantttitle{J}{1} \gantttitle{F}{1} \gantttitle{M}{1} \gantttitle{A}{1} \gantttitle{M}{1} \gantttitle{J}{1} \gantttitle{J}{1} \gantttitle{A}{1} \gantttitle{S}{1} \gantttitle{O}{1} \gantttitle{N}{1} \gantttitle{D}{1} \gantttitle{J}{1} \gantttitle{F}{1} \gantttitle{M}{1} \gantttitle{A}{1} \\

\ganttgroup{1. Literature review and review paper}{1}{10} \\
\ganttbar[bar/.append style={fill=gLit}]{Paper search and collection}{1}{6} \\
\ganttbar[bar/.append style={fill=gLit}]{Systematic reading and screening}{2}{8} \\
\ganttbar[bar/.append style={fill=gLit}]{Synthesis and gap analysis}{4}{9} \\
\ganttbar[bar/.append style={fill=gLit}]{Review paper draft}{7}{10} \\
\ganttmilestone{Review paper submitted}{10} \\

\ganttgroup{2. Research proposal}{3}{5} \\
\ganttbar[bar/.append style={fill=gProp}]{Proposal writing}{3}{5} \\
\ganttmilestone{Proposal defense}{5} \\

\ganttgroup{3. Programming and tooling}{2}{12} \\
\ganttbar[bar/.append style={fill=gProg}]{Python: fundamentals and ML libraries}{2}{7} \\
\ganttbar[bar/.append style={fill=gProg}]{MATLAB: toolboxes and sensor fusion}{5}{10} \\
\ganttbar[bar/.append style={fill=gProg}]{C++: image and point-cloud processing}{8}{12} \\

\ganttgroup{4. Dataset and preprocessing}{6}{16} \\
\ganttbar[bar/.append style={fill=gData}]{Data acquisition (RGB, thermal, LiDAR)}{6}{13} \\
\ganttbar[bar/.append style={fill=gData}]{Calibration and registration}{9}{14} \\
\ganttbar[bar/.append style={fill=gData}]{Labeling and annotation}{11}{15} \\
\ganttbar[bar/.append style={fill=gData}]{Reference depth measurement}{11}{15} \\
\ganttbar[bar/.append style={fill=gData}]{Multimodal preprocessing pipeline}{13}{16} \\
\ganttbar[bar/.append style={fill=gData}]{Registration accuracy validation}{14}{16} \\

\ganttgroup{5. Model development and evaluation}{13}{23} \\
\ganttbar[bar/.append style={fill=gModel}]{Baseline model (single modality)}{13}{16} \\
\ganttbar[bar/.append style={fill=gModel}]{Fused model (RGB, thermal and LiDAR)}{15}{19} \\
\ganttbar[bar/.append style={fill=gModel}]{Depth regression head}{16}{20} \\
\ganttbar[bar/.append style={fill=gModel}]{Validation and benchmarking}{21}{23} \\
\ganttbar[bar/.append style={fill=gModel}]{Ablation study and optimization}{20}{22} \\
\ganttmilestone{Model finalized}{23} \\

\ganttgroup{6. Thesis writing and dissemination}{17}{24} \\
\ganttbar[bar/.append style={fill=gThesis}]{Thesis chapters (draft)}{17}{23} \\
\ganttbar[bar/.append style={fill=gThesis}]{Journal paper preparation}{19}{23} \\
\ganttmilestone{Thesis submitted and viva}{24}
\end{ganttchart}

\vspace{0.6\baselineskip}
{\footnotesize
\begin{tabular}{@{}llll@{}}
\gkey{gPhase}~Phase span & \gkey{gLit}~Literature and review paper & \gkey{gProp}~Research proposal & \gkey{gProg}~Programming and tooling \\[2pt]
\gdia~Milestone & \gkey{gData}~Dataset and preprocessing & \gkey{gModel}~Model development & \gkey{gThesis}~Thesis and dissemination \\
\end{tabular}}

\caption{Study plan, May 2026 to April 2028. Each phase span marks the full
extent of a workstream; the bars beneath it are its constituent activities.}
\label{fig:gantt}
\end{figure}
\end{landscape}


\section{Preliminary Results}
\label{sec:preliminary-results}

A pilot experiment was conducted to exercise the ablation protocol of
Section~\ref{sec:ablation-protocol} before the trimodal dataset exists. Its
purpose is to establish that the protocol runs end to end and produces an
interpretable verdict --- not to produce a corrosion result, which it cannot,
for the reason given below.

\subsection{Rationale and scope}

Objective~3 claims that the contribution of each modality will be established
empirically rather than assumed. That claim is only worth making if the
machinery that would establish it exists and has been shown to work. Since no
public dataset pairs registered RGB and thermal imagery of corroded steel ---
the absence that motivates this research --- the pilot runs on LLVIP
\citep{Jia2021}, a set of 15{,}488 visible and long-wave thermal image pairs
registered at capture. Registration is therefore not a confounding variable, and
what the experiment isolates is the contribution of the channels themselves.

The pilot exercises the first two rungs of the Objective~3 ladder. It makes no
claim about corrosion, and no claim about LiDAR. Consistent with the sensing
requirement stated in Section~\ref{sec:research-design}, no channel in the pilot
is synthesised from another: the thermal plane is measured by a thermal camera.

\subsection{Design}

Four input encodings were compared. Each is a three-channel image, so the
backbone, the input stem and the parameter count are identical across arms and
any difference in the metric is attributable to the information in the channels
rather than to a change in model capacity.

\begin{table}[htbp]
\centering
\caption{Input encodings compared in the pilot. All four are three-channel, so
model capacity is held constant across arms.}
\label{tab:prelim-arms}
\begin{tabularx}{\textwidth}{@{}l l X@{}}
\toprule
\textbf{Arm} & \textbf{Channels} & \textbf{What it isolates} \\
\midrule
RGB & R, G, B & the single-modality appearance baseline \\
CbCr & Cb, Cr, constant & colour-space effect with no additional sensor \\
Thermal & T, T, T & the thermal channel alone \\
CbCr--IR & Cb, Cr, T & fusion of the two sensors \\
\bottomrule
\end{tabularx}
\end{table}

The fused encoding is the construction of \citet{Lim2021}, in which a
two-channel chroma image is stacked with a one-channel thermal image before any
feature extraction. Reusing it means the pilot tests the closest published
RGB-plus-thermal corrosion design rather than an invention of this research.

The CbCr arm exists to separate two effects that \citet{Lim2021} change
together. His fused input alters the colour space and adds a sensor in the same
step, so his comparison against RGB cannot attribute the gain to either.
Isolating the colour-space change is what makes the thermal contribution
measurable.

Detection uses a single-stage detector \citep{Ultralytics2023} under the fixed
recipe of Section~\ref{sec:ablation-protocol}: one architecture, one schedule,
one split drawn with a fixed seed and reused by every arm, colour-jitter
augmentation disabled throughout, and three seeds per arm. The primary metric is
mAP averaged over 0.5 to 0.95 on a held-out test split. The decision rule is the
one stated in Section~\ref{sec:ablation-protocol} and is applied mechanically to
whatever the runs produce.

\subsection{Result}

% Replace this table with prelim_ablation/results/table.tex once the runs finish.
\begin{table}[htbp]
\centering
\caption{Modality ablation on the pilot dataset. All arms share one
architecture, one training recipe and one data split; only the input channels
differ. Figures are the mean over three seeds, with the observed range in
parentheses.}
\label{tab:prelim-ablation}
\begin{tabular}{@{}lcccc@{}}
\toprule
\textbf{Input} & \textbf{mAP@0.5 (\%)} & \textbf{mAP@0.5:0.95 (\%)} & \textbf{Precision (\%)} & \textbf{Recall (\%)} \\
\midrule
RGB & --- & --- & --- & --- \\
CbCr & --- & --- & --- & --- \\
Thermal & --- & --- & --- & --- \\
CbCr--IR (fused) & --- & --- & --- & --- \\
\bottomrule
\end{tabular}
\end{table}

For reference, \citet{Lim2021} report the same comparison on steel bridge
corrosion, on thirty-six acquired image pairs. His figures are reproduced in
Table~\ref{tab:lim-ablation} because they are the only published modality
ablation on corrosion with both an RGB arm and a thermal arm, and because the
shape of that result --- rather than its magnitude --- is what the pilot tests
for.

\begin{table}[htbp]
\centering
\caption{Modality ablation on steel bridge corrosion reported by
\citet{Lim2021}. AP$_1$ is surface corrosion, AP$_2$ subsurface. The fused input
exceeds the best appearance-only arm on surface corrosion and the best
thermal-only arm on subsurface corrosion, so the gain is not a trade between the
two tasks.}
\label{tab:lim-ablation}
\begin{tabular}{@{}lccc@{}}
\toprule
\textbf{Input} & \textbf{AP$_1$ (\%)} & \textbf{AP$_2$ (\%)} & \textbf{mAP (\%)} \\
\midrule
\textbf{CbCr--IR (fused)} & \textbf{88.84} & \textbf{88.59} & \textbf{88.72} \\
RGB & 78.97 & 11.54 & 45.26 \\
CbCr & 83.13 & 3.45 & 43.29 \\
Thermal (lock-in amplitude) & 55.52 & 87.08 & 71.30 \\
Thermal (three raw frames) & 51.75 & 79.23 & 65.49 \\
\bottomrule
\end{tabular}
\end{table}

Three things in Table~\ref{tab:lim-ablation} bear directly on this research.
The appearance channels are almost blind to subsurface corrosion, at 11.54\% and
3.45\% average precision, which is the sensory argument for a thermal channel
stated as a measurement rather than as an expectation. The thermal channel alone
is the weaker detector of surface corrosion, at 55.52\%, because structural
features such as bolts and welded joints produce false alarms --- so neither
channel dominates and the case for fusion is symmetric. And the fused arm
exceeds each single-modality arm on that arm's own strength, which is the
outcome an examiner would reasonably doubt: fusion here is not a compromise
between two partial detectors.

The comparison carries two caveats. The result was obtained under active
excitation by halogen lamps at one metre, so it does not establish that a
passive thermal channel would behave the same way; and the reported figures
come from a twenty-image test set, half of which are rotations of the other
half, so the effective test set is ten independent images.

\subsection{What the pilot establishes, and what it does not}

It establishes that the ablation protocol of
Section~\ref{sec:ablation-protocol} is executable: the encodings, the fixed
recipe, the seed repeats, the metric collection and the mechanical application
of the decision rule all run, and the pipeline that produces them transfers
unchanged to the trimodal experiment. Adding LiDAR adds arms and a geometric
channel to the encoding; it does not change the protocol.

It does not establish anything about corrosion. The pilot's domain is not this
research's domain, its target class is not corrosion, and its registration was
solved by the dataset's authors rather than by this research. Those three gaps
are what Objectives~1 and~2 exist to close, and the pilot is evidence that the
instrument for closing them is built and working, not that the answer is already
known.
